\section{Preliminary Research Directions}
\label{sec:directions}

The initial survey suggests that the chiplet discussion should not be reduced to the observation that a processor contains several pieces of silicon. The important architectural questions concern \textbf{what is partitioned, why the partition is made, and what costs are introduced by moving communication across die boundaries}.

The chiplet portion of the paper is therefore expected to examine subjects such as:

\begin{itemize}
  \item Core Complex Dies and I/O Dies;
  \item Infinity Fabric and die-to-die communication;
  \item latency and bandwidth;
  \item process-node specialization;
  \item manufacturing and die-size considerations;
  \item scalability with core count;
  \item interconnect power; and
  \item cache organization and coherence.
\end{itemize}

The 3D V-Cache portion has an especially direct connection to Chapter 11. A hierarchical memory system attempts to satisfy a large fraction of memory accesses in relatively small, fast caches. Increasing last-level-cache capacity can reduce expensive accesses to lower levels for workloads that benefit from the additional working-set capacity. However, conventional cache expansion also involves area, latency, power, and implementation costs.

3D V-Cache is therefore interesting as an architectural case study because it modifies part of this trade-off: additional SRAM capacity can be integrated vertically through a stacked cache die.

This leads to several questions for later study:

\begin{itemize}
  \item Under what workload conditions does additional L3 capacity provide substantial performance improvement?
  \item How should increased cache capacity be evaluated against access latency and clock frequency?
  \item Why are some workloads more cache-sensitive than others?
  \item What thermal and power constraints arise when cache is vertically stacked?
  \item How does the cache hierarchy interact with the chiplet organization of the processor?
  \item To what extent are observed performance gains attributable to cache, core microarchitecture, frequency, memory bandwidth, or other factors?
\end{itemize}
